\documentclass[10pt,a4paper]{amsart}
\usepackage{amsmath,amssymb,amsthm,graphicx}
\usepackage{verbatim}
\newenvironment{reference}{\comment}{\endcomment}
%\newenvironment{reference}{}{}
\newenvironment{slogan}{\comment}{\endcomment}
\newenvironment{history}{\comment}{\endcomment}

% For commutative diagrams we use Xy-pic
\usepackage[all]{xy}

% We use 2cell for 2-commutative diagrams.
\xyoption{2cell}
% Theorem environments.
%
\theoremstyle{plain}
\newtheorem{theorem}[subsection]{Theorem}
\newtheorem{proposition}[subsection]{Proposition}
\newtheorem{lemma}[subsection]{Lemma}

\theoremstyle{definition}
\newtheorem{definition}[subsection]{Definition}
\newtheorem{example}[subsection]{Example}
\newtheorem{exercise}[subsection]{Exercise}
\newtheorem{situation}[subsection]{Situation}

\theoremstyle{remark}
\newtheorem{remark}[subsection]{Remark}
\newtheorem{remarks}[subsection]{Remarks}

\numberwithin{equation}{subsection}

% Macros
%
\def\lim{\mathop{\mathrm{lim}}\nolimits}
\def\colim{\mathop{\mathrm{colim}}\nolimits}
\def\Spec{\mathop{\mathrm{Spec}}}
\def\Hom{\mathop{\mathrm{Hom}}\nolimits}
\def\Ext{\mathop{\mathrm{Ext}}\nolimits}
\def\SheafHom{\mathop{\mathcal{H}\!\mathit{om}}\nolimits}
\def\SheafExt{\mathop{\mathcal{E}\!\mathit{xt}}\nolimits}
\def\Sch{\mathit{Sch}}
\def\Mor{\mathop{\mathrm{Mor}}\nolimits}
\def\Ob{\mathop{\mathrm{Ob}}\nolimits}
\def\Sh{\mathop{\mathit{Sh}}\nolimits}
\def\NL{\mathop{N\!L}\nolimits}
\def\CH{\mathop{\mathrm{CH}}\nolimits}
\def\proetale{{pro\text{-}\acute{e}tale}}
\def\etale{{\acute{e}tale}}
\def\QCoh{\mathit{QCoh}}
\def\Ker{\mathop{\mathrm{Ker}}}
\def\Im{\mathop{\mathrm{Im}}}
\def\Coker{\mathop{\mathrm{Coker}}}
\def\Coim{\mathop{\mathrm{Coim}}}

% Boxtimes
%
\DeclareMathSymbol{\boxtimes}{\mathbin}{AMSa}{"02}

%
% Macros for moduli stacks/spaces
%
\def\QCohstack{\mathcal{QC}\!\mathit{oh}}
\def\Cohstack{\mathcal{C}\!\mathit{oh}}
\def\Spacesstack{\mathcal{S}\!\mathit{paces}}
\def\Quotfunctor{\mathrm{Quot}}
\def\Hilbfunctor{\mathrm{Hilb}}
\def\Curvesstack{\mathcal{C}\!\mathit{urves}}
\def\Polarizedstack{\mathcal{P}\!\mathit{olarized}}
\def\Complexesstack{\mathcal{C}\!\mathit{omplexes}}
% \Pic is the operator that assigns to X its picard group, usage \Pic(X)
% \Picardstack_{X/B} denotes the Picard stack of X over B
% \Picardfunctor_{X/B} denotes the Picard functor of X over B
\def\Pic{\mathop{\mathrm{Pic}}\nolimits}
\def\Picardstack{\mathcal{P}\!\mathit{ic}}
\def\Picardfunctor{\mathrm{Pic}}
\def\Deformationcategory{\mathcal{D}\!\mathit{ef}}
\usepackage{cleveref}
\crefname{originalreference}{original reference}{original references}
\Crefname{originalreference}{Original reference}{Original references}
\makeatletter
\newlabel{lemma-derived-quasi-coherent-small-etale-site}{{071Q}{source}{}{sourcepointer-1}{}}
\newlabel{lemma-derived-quasi-coherent-small-etale-site@cref}{{[originalreference][0][]071Q}{[1][0][]source}{}{}{}}
\makeatother
\crefname{originalreference}{original reference}{original references}
\Crefname{originalreference}{Original reference}{Original references}
\makeatletter
\newlabel{lemma-descend-homomorphisms}{{09RE}{source}{}{sourcepointer-2}{}}
\newlabel{lemma-descend-homomorphisms@cref}{{[originalreference][0][]09RE}{[1][0][]source}{}{}{}}
\makeatother
\crefname{originalreference}{original reference}{original references}
\Crefname{originalreference}{Original reference}{Original references}
\makeatletter
\newlabel{lemma-descend-perfect}{{08HG}{source}{}{sourcepointer-3}{}}
\newlabel{lemma-descend-perfect@cref}{{[originalreference][0][]08HG}{[1][0][]source}{}{}{}}
\makeatother
\crefname{originalreference}{original reference}{original references}
\Crefname{originalreference}{Original reference}{Original references}
\makeatletter
\newlabel{lemma-glue}{{01JB}{source}{}{sourcepointer-4}{}}
\newlabel{lemma-glue@cref}{{[originalreference][0][]01JB}{[1][0][]source}{}{}{}}
\makeatother
\crefname{originalreference}{original reference}{original references}
\Crefname{originalreference}{Original reference}{Original references}
\makeatletter
\newlabel{lemma-induction-principle}{{08DR}{source}{}{sourcepointer-5}{}}
\newlabel{lemma-induction-principle@cref}{{[originalreference][0][]08DR}{[1][0][]source}{}{}{}}
\makeatother
\crefname{originallemma}{original lemma}{original lemmas}
\Crefname{originallemma}{Original lemma}{Original lemmas}
\makeatletter
\newlabel{perfect-lemma-descend-perfect}{{09RF}{source}{}{sourcepointer-6}{}}
\newlabel{perfect-lemma-descend-perfect@cref}{{[originallemma][0][]09RF}{[1][0][]source}{}{}{}}
\makeatother
\crefname{originalreference}{original reference}{original references}
\Crefname{originalreference}{Original reference}{Original references}
\makeatletter
\newlabel{situation-descent}{{084R}{source}{}{sourcepointer-7}{}}
\newlabel{situation-descent@cref}{{[originalreference][0][]084R}{[1][0][]source}{}{}{}}
\makeatother
\title{Perfect objects descend along limit transformations}
\author{The Stacks Project Authors}
\date{stacks-project:spaces-perfect.tex:local-snapshot-20261012; GNU FDL 1.2 or later}
\begin{document}
\maketitle
\noindent\textit{Editorial note.} All delivered mathematical text is the author's own, copied from the pinned Stacks Project chapter file spaces-perfect.tex; no statement or proof step was rewritten, reordered or removed.\par\medskip
\noindent\textit{Editorial note.} The counted proof is the author's complete argument: the source proves essential surjectivity by induction using the gluing lemma on a distinguished square and discharges the affine base case by the schemes result; the argument is complete\par\medskip
\noindent\textit{Editorial note.} The label lemma-perfect-on-limit names the counted statement, and the proof environment immediately following it in the source is delivered verbatim as its proof.\par\medskip
\noindent\textit{Editorial note.} The counted statement is carried by the author's own lemma environment in the chapter source.\par\medskip
\noindent\textit{Allowed standard prerequisites (editorial).}\begin{itemize}
\item the descent situation of the chapter: a limit of schemes or algebraic spaces with transition maps and X the limit
\item the induction principle for spaces etale over the limit
\item full faithfulness of the colimit functor on perfect objects (descend-homomorphisms lemma)
\item the gluing lemma for objects with maps over a distinguished square covering
\item perfect objects of the derived category
\end{itemize}\medskip
\ifdefined\hypertarget\hypertarget{sourcepointer-1}{}\fi
\noindent\textit{Source reference.} Original reference 071Q: sources/stacks-spaces-perfect/pinned-original-source.tex, source lines 402--402. Prerequisite pointer invoked at proof line(s) 6286; it names a standard fact of the same chapter that the counted argument uses, and every step of the counted proof built on it is retained verbatim in the delivered body.\par\medskip
\ifdefined\hypertarget\hypertarget{sourcepointer-2}{}\fi
\noindent\textit{Source reference.} Original reference 09RE: sources/stacks-spaces-perfect/pinned-original-source.tex, source lines 6184--6184. Prerequisite pointer invoked at proof line(s) 6249; it names a standard fact of the same chapter that the counted argument uses, and every step of the counted proof built on it is retained verbatim in the delivered body.\par\medskip
\ifdefined\hypertarget\hypertarget{sourcepointer-3}{}\fi
\noindent\textit{Source reference.} Original reference 08HG: sources/stacks-spaces-perfect/pinned-original-source.tex, source lines 3288--3288. Prerequisite pointer invoked at proof line(s) 6287; it names a standard fact of the same chapter that the counted argument uses, and every step of the counted proof built on it is retained verbatim in the delivered body.\par\medskip
\ifdefined\hypertarget\hypertarget{sourcepointer-4}{}\fi
\noindent\textit{Source reference.} Original reference 01JB: sources/stacks-spaces-perfect/pinned-original-source.tex, source lines 2206--2206. Prerequisite pointer invoked at proof line(s) 6274; it names a standard fact of the same chapter that the counted argument uses, and every step of the counted proof built on it is retained verbatim in the delivered body.\par\medskip
\ifdefined\hypertarget\hypertarget{sourcepointer-5}{}\fi
\noindent\textit{Source reference.} Original reference 08DR: sources/stacks-spaces-perfect/pinned-original-source.tex, source lines 1563--1563. Prerequisite pointer invoked at proof line(s) 6247; it names a standard fact of the same chapter that the counted argument uses, and every step of the counted proof built on it is retained verbatim in the delivered body.\par\medskip
\ifdefined\hypertarget\hypertarget{sourcepointer-6}{}\fi
\noindent\textit{Source reference.} Original lemma 09RF: sources/stacks-chow/tags.txt, source lines 12318--12318. Prerequisite pointer invoked at proof line(s) 6284; it names a standard fact of a sibling chapter that the counted argument uses, and every step of the counted proof built on it is retained verbatim in the delivered body.\par\medskip
\ifdefined\hypertarget\hypertarget{sourcepointer-7}{}\fi
\noindent\textit{Source reference.} Original reference 084R: sources/stacks-spaces-perfect/pinned-original-source.tex, source lines 6175--6175. Prerequisite pointer invoked at proof line(s) 6230; it names a standard fact of the same chapter that the counted argument uses, and every step of the counted proof built on it is retained verbatim in the delivered body.\par\medskip
\begin{lemma}
\label{lemma-perfect-on-limit}
In Situation \ref{situation-descent} the category of perfect
objects of $D(\mathcal{O}_X)$ is the colimit of the categories
of perfect objects of $D(\mathcal{O}_{X_i})$.
\end{lemma}

\begin{proof}
For every quasi-compact and quasi-separated object $U_0$ of
$(X_0)_{spaces, \etale}$ consider the condition $P$ that
the functor
$$
\colim_{i \geq 0} D_{perf}(\mathcal{O}_{U_i})
\longrightarrow
D_{perf}(\mathcal{O}_U)
$$
is an equivalence where ${}_{perf}$ indicates the full subcategory of
perfect objects and where $U = X \times_{X_0} U_0$ and
$U_i = X_i \times_{X_0} U_0$. We will prove $P$ holds for every $U_0$
by the induction principle of Lemma \ref{lemma-induction-principle}.
First, we observe that we already know the functor is fully faithful
by Lemma \ref{lemma-descend-homomorphisms}. Thus it suffices to prove
essential surjectivity.

\medskip\noindent
We first check condition (2) of the induction principle. Thus suppose
that we have an elementary distinguished square
$(U_0 \subset X_0, V_0 \to X_0)$ and that $P$ holds for
$U_0$, $V_0$, and $U_0 \times_{X_0} V_0$. Let $E$ be a perfect object
of $D(\mathcal{O}_X)$. We can find $i \geq 0$ and $E_{U, i}$ perfect on $U_i$
and $E_{V, i}$ perfect on $V_i$ whose pullback to $U$ and $V$ are isomorphic
to $E|_U$ and $E|_V$. Denote
$$
a : E_{U, i} \to (R(X \to X_i)_*E)|_{U_i}
\quad\text{and}\quad
b : E_{V, i} \to (R(X \to X_i)_*E)|_{V_i}
$$
the maps adjoint to the isomorphisms $L(U \to U_i)^*E_{U, i} \to E|_U$
and $L(V \to V_i)^*E_{V, i} \to E|_V$. By fully faithfulness, after
increasing $i$, we can find an isomorphism
$c : E_{U, i}|_{U_i \times_{X_i} V_i} \to E_{V, i}|_{U_i \times_{X_i} V_i}$
which pulls back to the identifications 
$$
L(U \to U_i)^*E_{U, i}|_{U \times_X V} \to E|_{U \times_X V} \to
L(V \to V_i)^*E_{V, i}|_{U \times_X V}.
$$
Apply Lemma \ref{lemma-glue}
to get an object $E_i$ on $X_i$ and a map $d : E_i \to R(X \to X_i)_*E$
which restricts to the maps $a$ and $b$ over $U_i$ and $V_i$.
Then it is clear that $E_i$ is perfect and that
$d$ is adjoint to an isomorphism $L(X \to X_i)^*E_i \to E$.

\medskip\noindent
Finally, we check condition (1) of the induction principle, in other
words, we check the lemma holds when $X_0$ is affine.
This follows from the case of schemes, see
Derived Categories of Schemes, Lemma \ref{perfect-lemma-descend-perfect}.
To see this use the equivalence of
Lemma \ref{lemma-derived-quasi-coherent-small-etale-site}
and use the translation of Lemma \ref{lemma-descend-perfect}.
\end{proof}

\end{document}
